\documentclass[11pt,letterpaper]{article}
\usepackage[T1]{fontenc}
\usepackage[utf8]{inputenc}
\usepackage{lmodern,microtype}
\usepackage[margin=1in]{geometry}
\usepackage{amsmath,amssymb,amsthm,mathtools}
\usepackage{booktabs,array,longtable}
\usepackage{graphicx,xcolor}
\usepackage{enumitem}
\usepackage{fancyhdr}
\usepackage[most]{tcolorbox}
\usepackage{listings}
\usepackage{hyperref}
\definecolor{navy}{HTML}{17324D}
\definecolor{teal}{HTML}{087E8B}
\definecolor{paleblue}{HTML}{F0F5FA}
\hypersetup{colorlinks=true,linkcolor=navy,citecolor=teal,urlcolor=teal,
 pdftitle={Sharper precision for subexponential 3SAT},
 pdfauthor={ChatGPT research analysis},
 pdfsubject={Proved Hermite precision refinement and conditional structured-mode acceleration}}
\pagestyle{fancy}
\fancyhf{}
\fancyhead[L]{\small\color{navy}Sharper precision for subexponential 3SAT}
\fancyhead[R]{\small Research memorandum}
\fancyfoot[C]{\thepage}
\setlength{\headheight}{14pt}
\setlength{\emergencystretch}{2em}
\setlist{topsep=4pt,itemsep=3pt,leftmargin=*}
\allowdisplaybreaks[2]
\numberwithin{equation}{section}
\newtheorem{theorem}{Theorem}[section]
\newtheorem{lemma}[theorem]{Lemma}
\newtheorem{proposition}[theorem]{Proposition}
\newtheorem{corollary}[theorem]{Corollary}
\theoremstyle{definition}
\newtheorem{hypothesis}[theorem]{Hypothesis}
\newtheorem{definition}[theorem]{Definition}
\theoremstyle{remark}
\newtheorem{remark}[theorem]{Remark}
\newcommand{\poly}{\operatorname{poly}}
\newcommand{\bits}{\{0,1\}}
\newcommand{\vTwo}{v_2}
\newcommand{\ceil}[1]{\left\lceil #1\right\rceil}
\newcommand{\floor}[1]{\left\lfloor #1\right\rfloor}
\newcommand{\srccommit}{\texttt{42907c9e999cc00d6abe7a50bfc1023250edd1b9}}
\newtcolorbox{statusbox}{colback=paleblue,colframe=navy,boxrule=.6pt,
 arc=1pt,left=9pt,right=9pt,top=8pt,bottom=8pt,breakable}
\lstset{basicstyle=\ttfamily\small,breaklines=true,columns=fullflexible,
 frame=single,rulecolor=\color{black!20},showstringspaces=false}
\title{\color{navy}\bfseries Sharper precision for subexponential 3SAT\par
 \large A proved schedule refinement and a conditional route to a better logarithmic exponent}
\author{Research memorandum prepared for Vladimir Reshetnikov}
\date{October 26, 2026}
\begin{document}
\maketitle
\begin{abstract}
We analyze the October 25 revision of OpenAI's subexponential exact-counting
algorithm, using a pinned source revision. We obtain a concrete refinement of
its tapered Hermite interpolation: count the interpolation conditions exactly,
then charge the actual two-adic loss of each recovery multiplier. The resulting
working precision is at most $(6+o(1))MD/u$, compared with the source schedule's
$(24+o(1))MD/u$. In the standard colored-partition estimate, this halves the
leading constant in the exponential bound for mode enumeration. We also solve
the associated real-valued minimax taper problem exactly and show why changing
the taper alone cannot improve the logarithmic power.

For a stronger improvement, we give an exact compressed tensor representation
of the mode sum and formulate a precise structured-summation hypothesis.
Under that hypothesis, a retuned algorithm would count arbitrary fixed-width
CNF in $\poly(n,m)2^{O(n\log\log n/\log n)}$ time and polynomial space, with
$\poly(n,m)2^{O(n/\log n)}$ at fixed clause density. The required new summation
algorithm is not proved here. Exact-arithmetic implementations check the
proved precision identities and the compressed representation on finite
instances.
\end{abstract}

\begin{statusbox}
\textbf{What is established.} A smaller, rigorously sufficient precision
schedule; an improved explicit exponential upper-bound constant; an exact
optimal taper for the full loss functional; and an exact mode-network
reformulation.

\medskip
\textbf{What is still open in this work.} An improvement of the general
$O(n\log\log n/\sqrt{\log n})$ exponent to a smaller asymptotic order. The
$O(n\log\log n/\log n)$ rate below is a conditional target, not a claimed
new 3SAT theorem. The supplied Python programs are algebraic prototypes and
schedule routines, not a full subexponential SAT solver.
\end{statusbox}
\clearpage
\tableofcontents
\clearpage

\section{The baseline and the result being improved}
\label{sec:baseline}

The user's reference is OpenAI, \emph{The Exponential Time Hypothesis is
False}, dated October 17 and revised October 25, 2026~\cite{paper}.
This memorandum uses repository commit \srccommit. The main theorem and
quantitative appendix state, for every fixed clause width $k$,
\begin{align}
 T_{\rm general}(n,m)
  &=(n+m)^{O_k(1)}2^{O_k(n\log\log n/\sqrt{\log n})},
   \label{eq:baseline-general}\\
 T_{\rm density}(n,m)
  &=(n+m)^{O_{k,D}(1)}2^{O_{k,D}(n/\sqrt{\log n})}
       \quad(m\le Dn,\ D\text{ fixed}).\label{eq:baseline-sparse}
\end{align}
Both are exact-counting bounds, so testing whether the count is positive
decides 3SAT when $k=3$. Both use polynomial space.

The distinction between the manuscript and the formalized implementation
matters. The repository's scope document~\cite{scope} describes a Lean theorem
with the weaker exponent $O(n(\log\log n)^2/\sqrt{\log n})$, and explicitly
lists the sharper manuscript bounds separately. The formal solution contains
a compiled counting machine~\cite{leanfinal}; it is not merely an informal
algorithm sketch. We inspected that distinction but did not rebuild Lean or
rerun its proof comparison machinery. The new work here is an analysis of
the quantitative tensor method, not a modification of that Lean endpoint.

It would therefore be misleading to present removal of one factor of
$\log\log n$ from the older bound as a new discovery. We compare against
\eqref{eq:baseline-general} and \eqref{eq:baseline-sparse}.

\subsection{A precise meaning of the proved improvement}

The change proved in Sections~\ref{sec:linear}--\ref{sec:global} improves
an explicit constant in the exponent of the tensor algorithm. It leaves its
big-$O$ logarithmic rate unchanged. If two implementations are certified by
$\exp((A+o(1))f(s))$ and $\exp((B+o(1))f(s))$ with $B<A$, the second has a
strictly smaller stated upper bound. This does not prove that the first
implementation attains its bound, or that a measured runtime ratio is
$\exp((A-B)f(s))$.

Our refinement has exactly that status. The more ambitious proposal in
Section~\ref{sec:conditional} would change the rate itself, but depends on
a new algorithmic hypothesis.

\begin{center}
\small
\begin{tabular}{@{}p{.25\linewidth}p{.43\linewidth}p{.22\linewidth}@{}}
\toprule
Result & Quantitative statement & Status \\
\midrule
Exact precision schedule & $P'_t=p+B_i+(q_i-1)(a+1-i)$ & Proved here \\
Dominant precision & $24MD/u$ replaced by at most $6MD/u$, asymptotically & Proved here \\
Enumeration bound & $4\pi\sqrt6$ replaced by $2\pi\sqrt6$ in the normalized exponent & Proved here \\
Mode-network encoding & $O(M/\sqrt u+r^2)$ rank bits for fixed weight allocations & Exact reformulation \\
General 3SAT target & $\poly(n,m)2^{O(n\log\log n/\log n)}$ & Conditional \\
Fixed-density target & $\poly(n,m)2^{O(n/\log n)}$ & Conditional \\
\bottomrule
\end{tabular}
\end{center}

\section{Where the source spends exponential time}
\label{sec:bottleneck}

We use the source's contraction theorem and its correctness argument as
the starting point~\cite{contraction,reduction}. A bipartite Boolean tensor
network with $M$ wires has contraction
\[
 Z=\sum_{x\in\bits^{E}}\prod_v T_v(x|_{\delta(v)}).
\]
The source handles suitably bounded degrees, polynomial-size descriptions,
and ring entries in
\[
 (\mathbb Z/2^{P_0}\mathbb Z)[V]/(V^\ell),
 \qquad P_0,\ell=O(s),
\]
where $s$ is a root size parameter and $M\le Cs$.
An unramified coefficient extension supplies residue-distinct interpolation
points. All precision statements below apply coordinatewise in that same
extension and truncated polynomial ring.

At an active call, the construction uses packets of $L=2^a$ slots and
integer interpolation nodes
\[
 S=\{-L,-L+1,\ldots,L-1\},\qquad \alpha=2^{a+1}.
\]
For a retained nonzero node $t$, put
\[
 i=\vTwo(t),\qquad D_i=2^{a-i},\qquad 0\le i<u\le a.
\]
The deformed contraction $F(T)$ has degree at most
\begin{equation}
 R=M(2L-1).\label{eq:degree}
\end{equation}
At each node, crossing interactions are written in mode variables $Z$.
A mode has a positive integer weight $w_Z$, and its coefficient is divisible
by $2^{w_Z}$. There are at most $6M/D_i$ modes of any one weight. Thus only
tuples satisfying
\begin{equation}
 \sum_Zw_Z\beta_Z<P_t\label{eq:budget}
\end{equation}
can matter modulo $2^{P_t}$.

For a fixed tuple, the sum factors into child contractions:
\begin{equation}
 F(t+\alpha\zeta)
   =\sum_{\beta:\,\sum w_Z\beta_Z<P_t}
                \prod_{g=1}^{r} Z_{g,\beta}\pmod {2^{P_t}}.
 \label{eq:mode-identity}
\end{equation}
Each mode has one source group and one target group. With the source choice
$r=64$, every child has at most $3M/32<M/8$ wires.

The source prescribes $h_{\rm old}=\ceil{3M/u}$, multiplicities
$q_i=h_{\rm old}(u-i)$ for $i<u$, and
\begin{equation}
 P_t^{\rm old}=p+8h_{\rm old}D_i+8M(a+1).
 \label{eq:old-precision}
\end{equation}
In its active regime this is $(24+o(1))MD_i/u$. Colored partitions then give
$\exp(O(M/\sqrt u))$ tuples. Near the root, $u=\Theta(\log s)$, producing the
square root of a logarithm in the denominator.

For arbitrary clause density, the reduction first produces sparse polynomial
queries. A retained incidence monomial has degree $K=O(n)$ among
$N=n(\log n)^{O(1)}$ possible incidences. Binary ranks for its incidence
subsets require
\begin{equation}
 O\!\left(\log\binom{N+K}{K}\right)=O(n\log\log n)
 \label{eq:dense-entropy}
\end{equation}
wires. Contracting a network at that size gives
\eqref{eq:baseline-general}. This explains the two distinct losses:
the square root comes from mode enumeration, while the extra $\log\log n$
comes from the reduction's representation~\cite{reduction}.

\section{First refinement: the exact linear taper}
\label{sec:linear}

The original multiplicity count deliberately uses a loose lower bound.
The exact count immediately improves its parameter.

\begin{lemma}[Exact number of interpolation conditions]
\label{lem:exact-jets}
Let $a\ge u\ge1$, $L=2^a$, and assign multiplicity
$q_i=h(u-i)$ to every node of valuation $i<u$, and zero otherwise.
The total number of Hermite conditions is
\begin{equation}
 Q=h A_{a,u},\qquad
 A_{a,u}=L\left(2u-2+2^{1-u}\right).
 \label{eq:jets}
\end{equation}
In particular $A_{a,u}$ is an integer.
\end{lemma}
\begin{proof}
For $0\le i<a$, exactly $L/2^i$ nonzero nodes in $S$ have valuation $i$.
Since $i<u\le a$, the required sum is
\[
 Q=Lh\sum_{i=0}^{u-1}(u-i)2^{-i}.
\]
The finite geometric sum and its derivative give
$\sum_{i=0}^{u-1}(u-i)2^{-i}=2u-2+2^{1-u}$.
\end{proof}

\begin{theorem}[Minimal sufficient integer slope]
\label{thm:linear}
For the degree bound \eqref{eq:degree}, the smallest positive integer slope
within this linear-taper family that guarantees ordinary Hermite recovery is
\begin{equation}
 \boxed{h_{\min}=\ceil{\frac{M(2L-1)+1}{L(2u-2+2^{1-u})}}.}
 \label{eq:hmin}
\end{equation}
If $u\to\infty$ and $M/u\to\infty$ in the source's active regime, then
$h_{\min}=(1+o(1))M/u$.
\end{theorem}
\begin{proof}
A polynomial of degree at most $R$ is uniquely determined by $R+1$
Hermite conditions at distinct nodes. Equation~\eqref{eq:jets} therefore
requires $hA_{a,u}\ge R+1$, giving \eqref{eq:hmin} and its minimality.
The denominator is $2L(u-1+2^{-u})$. Dividing the numerator by this
expression, and accounting for an additive rounding error below one,
gives the stated asymptotic formula.
\end{proof}

All retained valuation levels remain active. The packets, group labels,
interpolation nodes, local mode expansion, and child topology therefore need
no change. If $M\ge u$, the odd nodes alone imply $A_{a,u}\ge Lu$, and
$R+1\le2ML$. Hence $h_{\min}\le\ceil{2M/u}\le3M/u$ and
$q_i\le3M$. These bounds are at least as strong as the auxiliary
$h\le4M/u$, $q_i\le4M$ estimates used in the source.

This first change reduces the slope by a factor tending to three.
The second change avoids unnecessary slack in the precision charged to
each retained node.

\section{Second refinement: charge the exact recovery loss}
\label{sec:precision}

Extend $q_i$ by zero for $i\ge u$. For a node of valuation $i<u$, define
the nonnegative integer
\begin{equation}
 B_i=\sum_{j=1}^{a-i}j\,2^{a-i-j}(q_i-q_{i+j})
    =D_i\sum_{j=1}^{a-i}\frac{j}{2^j}(q_i-q_{i+j}).
 \label{eq:Bi}
\end{equation}
Nonnegativity holds for a nonincreasing taper, including both the linear
taper and the rounded full-loss optimum discussed below.

For clarity, introduce the Hermite cardinal factor
\[
 E_t(T)=\prod_{\substack{v\in S\setminus\{0,t\}}}
             \left(\frac{T-v}{t-v}\right)^{q_{\vTwo(v)}}.
\]
Zero multiplicities contribute a factor one. Counting the valuations of
differences in the complete residue grid gives the exact source identity
\begin{equation}
 \vTwo(E_t(0))=-B_i.\label{eq:E-valuation}
\end{equation}
This is the equality preceding the source's coarse bound on Hermite loss;
our change uses that equality directly~\cite{contraction}.

\begin{theorem}[Exact sufficient working precision]
\label{thm:precision}
Suppose $\deg F\le R$ and $F(0)$ is integral in the source ring.
At every retained node $t$, suppose $F(t+\alpha z)$ is integral
coefficientwise and its first $q_i$ coefficients are supplied modulo
$2^{P'_t}$, where
\begin{equation}
 \boxed{P'_t=p+B_i+(q_i-1)(a+1-i).}\label{eq:new-precision}
\end{equation}
For a nonnegative nonincreasing taper with enough Hermite conditions,
these jets determine $F(0)$ modulo $2^p$.
\end{theorem}
\begin{proof}
Every factor of $E_t(t+\alpha z)$ is a power of
$1+\alpha z/(t-v)$. Since $\vTwo(t-v)\le a$, its coefficients are integral
and its constant term is one. Its reciprocal truncated below $q_i$ is
therefore integral. Forming
\[
 G_t(z)=\left[F(t+\alpha z)/E_t(t+\alpha z)\right]_{<q_i}
\]
loses no precision. Hermite recovery evaluates this polynomial at
$-t/\alpha$ and multiplies by $E_t(0)$, then sums over $t$.

The multiplier of the coefficient of $z^j$, $0\le j<q_i$, has valuation
\[
 \vTwo\!\left(E_t(0)(-t/\alpha)^j\right)
       =-B_i-j(a+1-i).
\]
Replacing that coefficient by an arbitrary lift of its residue introduces
an error with valuation at least
\[
 P'_t-B_i-j(a+1-i)
       =p+(q_i-1-j)(a+1-i)\ge p.
\]
Every error lies in $2^p$ times the integral coefficient ring. Addition
does not reduce this precision. The exact Hermite identity recovers
$F(0)$, so the approximate calculation gives its required residue.
The argument holds separately in every coefficient coordinate, including
the truncated $V$ coordinate.
\end{proof}

\begin{remark}
This is a sufficient precision bound. It need not be minimal for a
particular polynomial with additional coefficient divisibility. For a
general nonmonotone taper, a safe nondecreasing-precision version replaces
the increment in \eqref{eq:new-precision} by its maximum with zero.
For the positive linear taper with $u\le a$, $B_i>0$, so precision strictly
increases, preserving the source's canonical-lifting convention.
\end{remark}

\subsection{Why the dominant constant becomes six}

Write $k_i=u-i$. The full normalized loss of a linear taper is
\begin{equation}
 \sum_{j\ge1}\frac{j}{2^j}(q_i-q_{i+j})
   =h\left(6-(k_i+3)2^{1-k_i}\right)\le6h.
 \label{eq:linear-loss}
\end{equation}
The finite sum in \eqref{eq:Bi} is no larger. Thus $B_i\le6h_{\min}D_i$.

The source uses cumulative packet exponents
\[
 u_d=\floor{\varepsilon\lambda/2^d},\quad
 a_d=2\sum_{j=0}^d u_j,\quad
 \lambda=\ceil{\log_2s},\quad
 \varepsilon=\frac1{128\cdot1024}.
\]
At every active node, $D_i\ge2^{a_{d-1}+u_d}$, while inherited precision is
at most $s\lambda^{O(1)}2^{a_{d-1}}$. The active $u_d$ is a positive
power of $\lambda$. Consequently, uniformly over active depths and retained
nodes,
\[
 p=o(MD_i/u),\qquad q_i(a+1-i)=o(MD_i/u).
\]
The smaller new increments preserve the inherited-precision estimate.
Combining these observations with Theorem~\ref{thm:linear} proves
\begin{equation}
 \boxed{P'_t\le(6+o(1))\frac{MD_i}{u},\qquad
 P_t^{\rm old}=(24+o(1))\frac{MD_i}{u}.}
 \label{eq:precision-comparison}
\end{equation}
Levels close to the end of the taper can do better than the upper constant
six. No change in the coefficient extension, exact arithmetic model, or
polynomial-space argument is required.

\section{The resulting explicit exponential bound}
\label{sec:global}

Let $N(P)$ be the number of mode tuples with weighted degree below $P$.
For any $\theta>0$, the multiplicity bound $6M/D_i$ gives
\begin{align}
 \log N(P)
 &\le\theta P+\frac{6M}{D_i}
                 \sum_{j\ge1}-\log(1-e^{-\theta j})\\
 &\le\theta P+\frac{\pi^2M}{D_i\theta}.
 \label{eq:entropy-bound}
\end{align}
Here and below $\log$ in an entropy calculation is natural. The second
line follows by expanding the logarithm:
\[
 \sum_{j\ge1}-\log(1-e^{-\theta j})
 =\sum_{h\ge1}\frac1{h(e^{\theta h}-1)}
 \le\frac1\theta\sum_{h\ge1}\frac1{h^2}
 =\frac{\pi^2}{6\theta}.
\]
Taking $\theta=\pi\sqrt{M/(D_iP)}$ yields
\begin{equation}
 \log N(P)\le2\pi\sqrt{MP/D_i}.\label{eq:optimized-entropy}
\end{equation}
Therefore \eqref{eq:precision-comparison} gives the explicit comparison
\begin{align}
 \log N(P_t^{\rm old})
   &\le(4\pi\sqrt6+o(1))M/\sqrt u,\\
 \log N(P'_t)
   &\le(2\pi\sqrt6+o(1))M/\sqrt u.
 \label{eq:mode-constant}
\end{align}

\begin{corollary}[An explicit bound for the complete contraction]
\label{cor:global}
Normalize the root so it has at most $s$ wires, and retain the source's
leaf threshold $H=\ceil{s/\floor{\sqrt\lambda}}$ and fixed $\varepsilon$.
With description length $\mathcal L$, the refined algorithm satisfies
\begin{equation}
 T\le(s+\mathcal L)^{O(1)}
 \exp\!\left[
 \left(\frac{2\pi\sqrt6}
 {\sqrt\varepsilon(1-\sqrt2/8)}+\log2+o(1)\right)
 \frac{s}{\sqrt\lambda}\right].
 \label{eq:global-explicit}
\end{equation}
The same calculation for the original precision prescription replaces
$2\pi\sqrt6$ by $4\pi\sqrt6$.
\end{corollary}
\begin{proof}
At depth $d$, $M_d\le s8^{-d}$ and
$u_d=(1+o(1))\varepsilon\lambda2^{-d}$ uniformly along active paths.
Thus the logarithms of the branching bounds sum to at most
\[
 \frac{2\pi\sqrt6+o(1)}{\sqrt\varepsilon}
 \frac{s}{\sqrt\lambda}
 \sum_{d\ge0}(\sqrt2/8)^d.
\]
Polynomial factors across $O(\log\lambda)$ levels contribute only
$O(\lambda\log\lambda)=o(s/\sqrt\lambda)$ to the logarithm.
Direct enumeration at a leaf contributes
$(\log2+o(1))s/\sqrt\lambda$. The allowed entry cost
$\exp(O(s^{9/10}))$ is also lower order. The fixed polynomial dependence
on description length is retained exactly as in the source analysis.
\end{proof}

This is a genuine refinement of the specified algorithm and its explicit
upper bound. Its asymptotic class remains
$\poly(s+\mathcal L)2^{O(s/\sqrt{\log s})}$. In particular, the manuscript's
general and fixed-density rates in Section~\ref{sec:baseline} remain the
strongest rates established here. The large fixed constants and eventual
parameter thresholds in the source also remain; this analysis is not a
practical performance claim.

\section{An exact optimal-taper theorem}
\label{sec:optimal}

The preceding improvement raises a natural question: can a more elaborate
taper reduce precision by an additional factor of $u$? Within the loss
functional used by the source, the answer is no. We can solve the relevant
real optimization problem exactly.

For a nonnegative real vector $q_0,\ldots,q_{u-1}$, set $q_i=0$ for
$i\ge u$, and define
\[
 \mathcal A_i(q)=\sum_{j\ge1}\frac{j}{2^j}(q_i-q_{i+j}),
 \qquad W(q)=\sum_{i=0}^{u-1}2^{-i}q_i.
\]

\begin{theorem}[Optimal full normalized loss]
\label{thm:optimal}
For $W>0$, among nonnegative real tapers satisfying $W(q)\ge W$,
\begin{equation}
 \min_q\max_{0\le i<u}\mathcal A_i(q)
   =\frac{W}{u/3+2/9-(2/9)4^{-u}}.
 \label{eq:optimal-value}
\end{equation}
Writing the right side as $\eta$, an optimizer is
\begin{equation}
 q_i=\eta r_{u-i},\qquad
 r_k=\frac{k}{6}+\frac5{18}+\frac29\,4^{-k}.
 \label{eq:optimal-profile}
\end{equation}
Every full loss of this profile equals $\eta$.
\end{theorem}
\begin{proof}
Since $\sum_{j\ge1}j2^{-j}=2$, the constraints
$\mathcal A_i(q)\le\eta$ are equivalent to
\[
 q_i\le\frac\eta2+\sum_{j\ge1}\frac{j}{2^{j+1}}q_{i+j}.
\]
Backward induction shows that the pointwise largest feasible vector
saturates each inequality. After dividing by $\eta$ and reversing the
index, its coefficients satisfy
\[
 r_k=\frac12+\sum_{j=1}^{k-1}\frac{j}{2^{j+1}}r_{k-j}.
\]
Their generating function is
\[
 \sum_{k\ge1}r_kz^k
   =\frac{z(2-z)^2}{2(1-z)^2(4-z)}.
\]
Partial fractions give \eqref{eq:optimal-profile}; substitution also
verifies it directly. Summing its weighted mass yields
\[
 \sum_{i=0}^{u-1}2^{-i}r_{u-i}
    =\frac u3+\frac29-\frac29\,4^{-u}.
\]
The maximal mass at loss $\eta$ is therefore $\eta$ times this quantity.
Choosing $\eta$ to attain mass $W$ proves both directions.
\end{proof}

Hermite completeness asks for $W=(R+1)/L=2M+o(M)$. Thus the optimal
full normalized loss is $(6+o(1))M/u$. The simple optimized linear taper
already reaches this leading constant.

\subsection{The finite grid and integer multiplicities}

Let $\mathcal A_i^{(a)}$ truncate the defining sum after $j=a-i$.
Its relation to the full loss is exact:
\begin{equation}
 \mathcal A_i(q)-\mathcal A_i^{(a)}(q)
    =q_i(a-i+2)2^{-(a-i)}.\label{eq:finite-tail}
\end{equation}
For $a\ge2u$ and $q_i=O(M)$, the right side is $o(M/u)$ uniformly as
$u\to\infty$. Hence, among bounded tapers in the source regime,
the finite-grid minimax loss has the same leading constant six.

An exact finite-grid optimizer can also be computed by a triangular
rational recurrence. Put
\[
 s_i=2-(a-i+2)2^{-(a-i)},\qquad
 r_i^{(a)}=\frac{1+\sum_{j=1}^{u-1-i}j2^{-j}r_{i+j}^{(a)}}{s_i},
\]
processing $i$ backwards. Scale this vector to mass $W$.
The same backward-induction proof shows optimality for the finite
normalized-loss problem. This optional finite optimizer need not be used
to obtain our leading-constant improvement.

For either real optimizer, rounding every coordinate upward preserves
Hermite completeness. Each loss increases by less than two: its diagonal
coefficient is at most two, while all off-diagonal coefficients are
nonpositive. Thus integer rounding preserves the leading constant.
For the rounded full profile, monotonicity follows from that of $r_k$;
the precision formula in Theorem~\ref{thm:precision} applies directly.

\subsection{A reproducible finite comparison}

For $M=1000$, $a=12$, $u=6$, and input precision $p=1001$, the supplied
exact schedule routine gives the following maximum over retained nodes:
\begin{center}
\begin{tabular}{@{}lr@{}}
\toprule
Schedule & Maximum working precision, in bits \\
\midrule
Published source prescription & $16{,}489{,}001$ \\
Minimal complete linear taper, exact loss & $4{,}684{,}588$ \\
Rounded full-loss optimum, exact loss & $3{,}681{,}886$ \\
Rounded finite-grid optimum, exact loss & $3{,}640{,}672$ \\
\bottomrule
\end{tabular}
\end{center}
These are exact integer schedule values, not timing measurements. The
example concerns the local recovery lemma, which does not require the
full solver's packet-packing inequalities. At small retained ranges, the
lower-order improvement from the optimal
profiles can be substantial. Upward rounding of a real optimum is not
claimed to solve the integer minimax problem exactly.

\subsection{What this rules out}

Within the retained-level architecture, a bounded taper satisfying
completeness must incur normalized loss $\Omega(M/u)$ at some active
level. At that level the Hermite multiplier loses
$\Omega(MD_i/u)$ bits. Therefore merely changing the taper cannot replace
that scale by $MD_i/u^2$.

This is a restriction on a particular precision accounting method. It is
not a lower bound on all algorithms. A different recovery identity,
additional divisibility specific to the generated polynomials, or
cancellation between recovery terms could evade its hypotheses.

\section{Why ordinary mode pruning does not improve the rate}
\label{sec:barriers}

Two valid refinements look promising at first. Since the expansion is a
product of $M$ affine row factors, it satisfies $\sum_Z\beta_Z\le M$.
Also, the true valuation of a mode coefficient can exceed its assigned
weight by a baseline shift $a-i$. Neither restriction alone changes
the worst-case counting scale.

\begin{proposition}[A feasible binary family]
\label{prop:mode-family}
Consider an abstract mode family with $c=\Theta(M/D)$ colors per weight
and budget $P=\Theta(MD/u)$, in a regime where
$D/\sqrt u\to\infty$. There are parameter choices of this order for which
at least $2^{\Omega(M/\sqrt u)}$ feasible tuples remain even after imposing
$\beta_Z\in\{0,1\}$ and $\sum_Z\beta_Z\le M$.
\end{proposition}
\begin{proof}
Choose $J=\floor{D/(C\sqrt u)}$ for a sufficiently large fixed $C$ and
select all colors at weights between $J$ and $2J$. There are
$K=\Theta(cJ)=\Theta(M/\sqrt u)$ such modes. Each of their $2^K$ subsets
has ordinary degree at most $K\le M$ and weight at most
$2JK=O(MD/u)$, below the budget after choosing $C$ appropriately.
\end{proof}

In the source's regime, the uniform baseline shift is negligible compared
with $J$; it does not remove this family. Balanced versions also satisfy
per-packet ordinary-degree caps. This proves a barrier to improving an
enumeration bound from those restrictions alone. It does not assert that
all coefficients in every actual network are nonzero.

Changing polynomial bases offers a related constant-level improvement,
but no established new rate: factorial divisibility obeys
$\vTwo(b!)=b-s_2(b)$, which is asymptotically linear in $b$. It does not
provide the superlinear weight growth needed to change the partition
exponent. The refinement proved above extracts the available taper
constant without claiming such a change.

\section{An exact compressed representation of the mode sum}
\label{sec:compressed}

The stronger proposal is to exploit the structure of
\eqref{eq:mode-identity} instead of attempting to enumerate fewer abstract
tuples. Its source--target incidence suggests an exact second tensor
network.

Partition the modes by ordered source--target pair $(g,h)$. Let
$\mathcal Z_{gh}$ be one such class. For an integer weight $w$, define
\[
 \mathcal C_{gh}(w)=
 \{(\beta_Z)_{Z\in\mathcal Z_{gh}}:
                       \sum_{Z\in\mathcal Z_{gh}}w_Z\beta_Z=w\},
 \qquad c_{gh}(w)=|\mathcal C_{gh}(w)|.
\]
Enumerate arrays $(w_{gh})$ with total below $P_t$, skipping arrays whose
classes are empty. For a fixed array, encode each class tuple by its
lexicographic rank, using
\[
 b_{gh}=\ceil{\log_2 c_{gh}(w_{gh})}
\]
bits, with zero bits when the count is one. Unused binary rank strings
give tensor entry zero.

The tensor at group $g$ decodes the ranks of all incident mode classes
and takes the value $Z_{g,\beta}$ from the source's child construction.
Each tuple is visible to precisely its source and target. Summing over
the rank bits therefore gives exactly the contribution of that weight
array. Summing over arrays gives \eqref{eq:mode-identity}.

\begin{proposition}[Size and exactness of the rank encoding]
\label{prop:rank}
For a fixed feasible array, the number of rank bits is
\begin{equation}
 S_\beta\le \log_2 N(P_t)+r(r-1)=O(M/\sqrt u+r^2).
 \label{eq:rank-budget}
\end{equation}
The encoding is counting preserving. Completion counts and rank decoding
use polynomial time and space in the original root parameters. The array
enumeration has at most $P_t^{O(r^2)}$ possibilities.
\end{proposition}
\begin{proof}
For the chosen array, its tuple count is the product of its class counts.
It is at most the total feasible tuple count $N(P_t)$. Taking logarithms
and allowing one rounding bit for each ordered pair proves
\eqref{eq:rank-budget}. Class decoding is bijective on valid ranks, and
invalid ranks contribute zero, so the contraction neither drops nor
duplicates a tuple.

For modes with weights $v_1,\ldots,v_b$, suffix completion counts satisfy
\[
 C(j,w)=\sum_{e=0}^{\floor{w/v_j}}C(j+1,w-ev_j),
\]
with terminal counts $C(b+1,0)=1$ and zero at other weights. These exact
counts support ranking and unranking by subtraction of successive
completion blocks. All mode counts and precision parameters are polynomial
in $s$, and their count bit lengths are polynomial as well. The number of
weight coordinates is at most $r(r-1)$, proving the array bound.
\end{proof}

The new graph may be made bipartite by subdividing every rank wire with
an equality tensor, at a factor-two edge cost. This is an exact
reformulation, not yet a faster contraction algorithm.

\subsection{Why applying the published theorem twice is invalid}

The compressed graph has only $r$ genuine group tensors. In a balanced
family their degrees are of order $S_\beta/r$. If $r$ is fixed or
polylogarithmic, this exceeds the published theorem's
$S_\beta^{9/10}$ maximum-degree requirement. Fast entry evaluation does
not remove that requirement.

The packing inequality makes the obstruction quantitative. Its array
parameter is $D_{\rm arr}=r_{\rm label}^{2a}$, and it requires
\[
 M_{\rm call}\ge2D_{\rm arr}^{\,2}\Delta.
\]
For outer degree $\Delta\asymp S_\beta/r$, it allows only
$u'=O(\log r)$ retained levels. Increasing the number of initial groups
to $r$ simultaneously restricts the initial range to
$u=O(\log M/\log r)$. The apparent two savings then cancel:
\[
 \frac{M}{\sqrt{\log M/\log r}}\,
       \frac1{\sqrt{\log r}}
   =\frac{M}{\sqrt{\log M}}.
\]
This is a diagnosis of this particular attempted proof, not a general
algorithmic lower bound.

An equality tensor can be split into a bounded-degree tree because it has
special structure. A generic high-degree tensor cannot be split that way
merely from an efficient procedure for querying its entries. Retaining the
original child networks does supply structure, but restores $\Theta(M)$
original wires and loses the compression. A successful new theorem must
use the actual generated representation more effectively.

\section{The proposed improvement in logarithmic rate}
\label{sec:conditional}

The exact target can now be stated without hiding the missing step.

\begin{hypothesis}[Structured-mode acceleration]
\label{hyp:acceleration}
For the actual packet tensors and canonical descriptions generated by the
source, there is a uniform deterministic procedure evaluating
\eqref{eq:mode-identity} using at most
\[
 s^{O(1)}\exp(O(M/u))
\]
recursive contractions of the same admissible type, each with at most
$M/8$ wires. It preserves the source's degree and precision bounds,
polynomial description lengths, and polynomial space. Its other work has
only the source's allowed polynomial factors and entry-evaluation factor
$\exp(O(s^{9/10}))$. The recursive query count is independent of the
original description length; that length enters as a fixed polynomial
per-call work factor.
\end{hypothesis}

This is an algorithmic hypothesis about structured expressions, not about
arbitrary unrelated summands supplied through a black-box oracle. The
explicit mode enumeration in the paper does not meet it. The rank
representation above is a possible starting point, but does not prove it.

\begin{theorem}[Conditional faster contraction]
\label{thm:conditional-contraction}
Under Hypothesis~\ref{hyp:acceleration}, the source's class of tensor
networks can be contracted in
\[
 (s+\mathcal L)^{O(1)}2^{O(s/\log s)}
\]
time and polynomial space.
\end{theorem}
\begin{proof}
Replace the leaf threshold by $H=\ceil{s/\lambda}$, retaining
$u_d=\floor{\varepsilon\lambda/2^d}$ and $a_d=2\sum_{j\le d}u_j$.
At an active call,
\[
 s/\lambda<M\le Cs8^{-d},\qquad
 d\le\tfrac13\log_2\lambda+O_C(1),\qquad
 u_d=\Omega_C(\lambda^{2/3}).
\]
Thus one may replace the source's auxiliary lower bound
$u\ge\lambda^{4/5}$ by $u\ge\lambda^{3/5}$, for sufficiently large $s$.
The cumulative bound $a_d\le4\varepsilon\lambda$ is unchanged.

In particular, $L,D_{\rm arr}\le s^{1/1000}$ still holds, and
\[
 2D_{\rm arr}^{\,2}\Delta=O(s^{9/10+2/1000})=o(s/\lambda),
 \qquad LD_{\rm arr}=o(s/\lambda),\qquad 2L\le\Delta.
\]
All packing conditions survive. The inherited precision is still
$s\lambda^{O(1)}2^{a_{d-1}}$, while
$D_i\ge2^{a_{d-1}+u_d}$. Since $2^{u_d}$ dominates every fixed power of
$\lambda$, the existing $P_t=O(MD_i/u_d)$ bound survives as well.

Under the hypothesis, the logarithmic branching contribution at depth $d$
is $O(M/u_d)=O((s/\lambda)4^{-d})$. Summing the geometric series gives
$O(s/\lambda)$. Polynomial overhead across all depths contributes only
$O(\lambda\log\lambda)$. A leaf costs
\[
 \exp(O(H)+O(s^{9/10}))=\exp(O(s/\lambda)).
\]
The assumed streaming evaluation and immutable polynomial descriptions
permit the same depth-first polynomial-space traversal.
\end{proof}

\begin{corollary}[Conditional fixed-density counting]
\label{cor:conditional-sparse}
For fixed $k$ and $D$, Hypothesis~\ref{hyp:acceleration} implies exact
counting of formulas with $m\le Dn$ in
\[
 (n+m)^{O_{k,D}(1)}2^{O_{k,D}(n/\log n)}
\]
time and polynomial space, using the source's $O_{k,D}(n)$-wire clause
network.
\end{corollary}

\subsection{Two additional changes needed for arbitrary density}

Changing only the contraction exponent is insufficient. The original
general-density proof contains two separate factors
$2^{O(n/\sqrt{\log n})}$, which would exceed the proposed faster target.
Both must be retuned.

Let $\lambda=\ceil{\log_2n}$ and $\nu=\ceil{\log_2\lambda}$. First set
the sparsification parameter to
\begin{equation}
 \sigma=1/\lambda.\label{eq:sigma}
\end{equation}
The source's disjoint sparsification lemma gives
$\poly(n,\lambda)2^{n/\lambda}$ tree nodes and occurrence bound
\[
 B=O_k((\lambda\log(2\lambda))^{k-1})=\lambda^{O_k(1)}.
\]
The initial parity checks, greedy candidate lengths, and sparse polynomial
queries therefore still have polylogarithmic sparsity.

Second, if $N_0$ is the number of test--variable incidences in a query,
use the high-incidence cutoff
\begin{equation}
 D_* =\max\{1,\ceil{N_0/n}\}\lambda.\label{eq:high-cutoff}
\end{equation}
At most $n/\lambda$ variables exceed it. Enumerating their assignments
costs at most $2^{n/\lambda}$. The remaining scopes and degrees still
carry only fixed polylogarithmic factors.

The incidence-monomial encoding continues to have $s=n\nu$ as a valid
size parameter. Its maximum degree is at most
$n^{3/4}\lambda^{O_k(1)}<(n\nu)^{9/10}$ for sufficiently large $n$,
and its entry cost remains $\exp(\lambda^{O_k(1)})$.
Theorem~\ref{thm:conditional-contraction} would give exponent
$O(s/\log s)=O(n\nu/\lambda)$. The two new branching exponents
$O(n/\lambda)$ and the degree-array exponent $O(\sqrt n\log n)$ are
absorbed. Candidate enumeration contributes only fixed powers of
$\lambda$ to the exponent.

\begin{corollary}[Conditional general-density counting]
\label{cor:conditional-general}
Under Hypothesis~\ref{hyp:acceleration}, for every fixed $k$ one obtains
\begin{equation}
 \boxed{(n+m)^{O_k(1)}
       2^{O_k(n\log\log n/\log n)}}\label{eq:conditional-target}
\end{equation}
time and polynomial space for exact counting. In particular, this would
give the same bound for deciding 3SAT.
\end{corollary}

The argument applies to every queried sparse system, including candidates
that still admit assignments violating the original formula. No unproved
assumption about the greedy procedure's intermediate outputs is needed.

\section{Research priorities and decision criteria}
\label{sec:research}

The immediate implementation task is well defined: use
\eqref{eq:hmin}, \eqref{eq:Bi}, and \eqref{eq:new-precision} in a
quantitative contraction implementation, retaining the existing algebraic
interfaces and canonical-lifting rules. The supplied schedule module
provides exact integer and rational routines for this part. A future Lean
formalization should first prove these local identities, then transport
the source contraction correctness argument; it should not simply edit
the unrelated older formalized complexity endpoint.

For a new logarithmic rate, the most focused next task is a factorization
theorem for the generated high-degree tensors in Section~\ref{sec:compressed}.
A candidate must do all of the following:
\begin{enumerate}
\item Preserve the exact weighted budget and the coefficient extraction at
      every target group, including polynomial factors in the truncated
      ring and creation-time representatives.
\item Use the tensor descriptions or their source--target structure, so
      that the exponentially many feasible tuples need not be queried
      separately.
\item Bound total time for computing the factorization, not just its
      number of rank bits. Recomputing child entries can dominate a small
      representation.
\item Avoid reintroducing $\Theta(M)$ uncompressed wires or violating the
      degree hypothesis at the next contraction level.
\item Give a polynomial-space streaming evaluation and a uniform bound
      valid at all active recursion depths.
\end{enumerate}

A representation satisfying these conditions with the query bound in
Hypothesis~\ref{hyp:acceleration} would be a substantive breakthrough.
Merely showing that the sum has a smaller entropy encoding does not meet
them. The finite prototype is useful for falsifying proposed identities
early; it cannot establish the required asymptotic representation theorem.

The taper optimum also supplies a practical stopping rule for this
research direction. Further optimization of the same retained-level loss
can improve lower-order terms and constants, but cannot generate an extra
factor $1/\sqrt{\log s}$ in the exponent. That gain must enter through
new algebraic summation, a different recovery scheme, or additional
structure not captured by the present loss functional.

\section{Artifacts and reproducibility}
\label{sec:artifacts}

The archive contains this LaTeX source and its compiled PDF; exact-arithmetic
Python routines; generated verification records; source provenance; and
supporting derivations. The Python checks use the standard library.
\begin{description}[style=nextline]
\item[\texttt{code/precision\_schedule.py}]
Computes source, optimized-linear, and rounded optimal schedules using
integers and exact rational arithmetic. It exposes the proved local
refinement without claiming to execute the entire SAT algorithm.
\item[\texttt{code/verify\_precision\_schedule.py}]
Checks schedule completeness, exact loss accounting, and finite recovery
properties; its JSON result records the executed tests.
\item[\texttt{code/mode\_network.py}]
Implements a small exact Cauchy deformation, mode-polynomial coefficient
extraction, completion-count ranking, and the compressed mode network.
\item[\texttt{code/verify\_mode\_network.py}]
Compares the direct $X,Y$ sum, global mode enumeration, and pairwise-rank
contraction. Its fixtures include signed interpolation nodes, zero and
nonzero shifts, invalid rank codes, and a nonzero quadratic mode term.
\item[\texttt{code/verify\_bounds.py}]
Independently checks finite Hermite valuations by rational products,
finite-tail identities, taper optima, and feasible-family witnesses.
\item[\texttt{results/}]
Contains machine-readable verification summaries and exponent-scale
comparisons. The latter compare mathematical functions, not runtime
measurements. The faster rate is labeled conditional.
\item[\texttt{notes/} and \texttt{source\_manifest.json}]
Contain detailed derivations, the conditional theorem's assumptions, and
pinned source identities. No unproved acceleration routine is supplied
under an implementation label.
\end{description}

The mode-representation verification includes 11 exact identity fixtures,
133,376 direct assignment pairs, 847 feasible global mode tuples, and
456 invalid binary rank strings assigned zero. All three evaluations
agree. Nine accompanying unit tests pass. These are checks of the
reformulation on finite examples, not a new verification of the full
source paper or its asymptotic theorem.

The independent bounds verifier includes 5,460 exhaustive small taper
profiles, 92,840 finite-tail coordinate checks, and 46 direct rational
Hermite multiplier checks. The JSON summaries give the full parameters
and exact results; subsequent reruns should be compared with those
records rather than with wall-clock timing.

The precision module passes seven unit tests and 40 exact Hermite
recovery fixtures, with 304 independent multiplier-valuation checks and
304 precision-boundary checks. The fixtures include genuinely rational
Cauchy-product polynomials as well as integer polynomials. Errors at the
permitted precision preserve the recovered residue; decreasing the
allowance by one bit can change it, as checked by full rational recovery.

To reproduce the checks and rebuild the report from the archive root:
\begin{lstlisting}
python3 code/verify_precision_schedule.py
python3 code/verify_mode_network.py
python3 code/verify_bounds.py
latexmk -pdf -interaction=nonstopmode -halt-on-error report.tex
\end{lstlisting}

\section{Assessment}

The established contribution is a concrete improvement to the
quantitative implementation's schedule: the exact interpolation condition
count and exact valuation charge reduce its dominant precision coefficient
from 24 to at most 6, and halve the corresponding colored-partition
upper-bound coefficient. The optimal-taper theorem shows that the improved
leading precision coefficient is already optimal within the full
retained-level loss model.

The desired improvement in logarithmic order remains a research proposal.
Its missing ingredient is isolated in Hypothesis~\ref{hyp:acceleration};
the required recursion and both density-reduction changes have been worked
out. The exact compressed representation supplies an object on which to
attack that ingredient, while the degree and enumeration barriers explain
why it is not a consequence of the published contraction theorem alone.

\clearpage
\begin{thebibliography}{9}
\bibitem{paper}
OpenAI. \emph{The Exponential Time Hypothesis is False}.
October 17, 2026, revised October 25, 2026.
\href{https://github.com/openai/math/blob/42907c9e999cc00d6abe7a50bfc1023250edd1b9/preprints/The-Exponential-Time-Hypothesis-is-False-October-25-2026/paper.pdf}{Pinned paper} and
\href{https://github.com/openai/math/blob/42907c9e999cc00d6abe7a50bfc1023250edd1b9/preprints/The-Exponential-Time-Hypothesis-is-False-October-25-2026/build/main_full_version.tex}{main LaTeX source}.

\bibitem{contraction}
OpenAI. Quantitative contraction appendix to~\cite{paper}.
\href{https://github.com/openai/math/blob/42907c9e999cc00d6abe7a50bfc1023250edd1b9/preprints/The-Exponential-Time-Hypothesis-is-False-October-25-2026/build/quantitative-contraction.tex}{Pinned source: quantitative-contraction.tex}.
Especially the tapered Hermite recovery, mode multiplicity, mode budget,
recursion identity, and final time analysis.

\bibitem{reduction}
OpenAI. Quantitative reduction appendix to~\cite{paper}.
\href{https://github.com/openai/math/blob/42907c9e999cc00d6abe7a50bfc1023250edd1b9/preprints/The-Exponential-Time-Hypothesis-is-False-October-25-2026/build/quantitative-reduction.tex}{Pinned source: quantitative-reduction.tex}.
Especially the Boolean contraction interface, fixed-density clause
network, incidence degree, wire budget, and arbitrary-density proof.

\bibitem{scope}
OpenAI. Formalization scope for \emph{The Exponential Time Hypothesis is
False}. \href{https://github.com/openai/math/blob/42907c9e999cc00d6abe7a50bfc1023250edd1b9/lean/docs/378.md}{Pinned scope document, lean/docs/378.md}.

\bibitem{leanfinal}
OpenAI. Formal counting-machine endpoint.
\href{https://github.com/openai/math/blob/42907c9e999cc00d6abe7a50bfc1023250edd1b9/lean/OAI/Algorithms/SubexponentialSAT/SqrtLogSpace/Final.lean}{Pinned source: SqrtLogSpace/Final.lean}.
Inspected for implementation scope; not rebuilt in this investigation.
\end{thebibliography}
\end{document}
